\documentclass{article}
\usepackage{hyperref}
\usepackage{xcolor}

\usepackage{fontspec}
\setmainfont[Path="V:/Fonts/Eng/",AutoFakeBold=3,AutoFakeSlant=0.25]{BASKVILL.TTF}
\setsansfont[Path="V:/Fonts/Eng/",AutoFakeBold=3,AutoFakeSlant=0.25]{POORICH.TTF}
\setmonofont[Path="V:/Fonts/Eng/",AutoFakeBold=3,AutoFakeSlant=0.25]{cour.ttf}

\usepackage{xeCJK}
\setCJKmainfont[Path="V:/Fonts/Zh/",AutoFakeSlant=0.2,AutoFakeBold=3]{STZHONGS.TTF}
\setCJKsansfont[Path="V:/Fonts/Zh/",AutoFakeSlant=0.2,AutoFakeBold=3]{STKAITI.ttf}
\setCJKmonofont[Path="V:/Fonts/Zh/",AutoFakeSlant=0.2,AutoFakeBold=3]{simfang.ttf}

\usepackage[most]{tcolorbox}

% 定义计数器
\newcounter{requirement}[subsection]
\renewcommand{\therequirement}{R-\arabic{section}.\arabic{subsection}.\arabic{requirement}}

% 定义环境
\NewDocumentEnvironment{requirement}{oo}{%
    \refstepcounter{requirement}%
    \par\noindent
    \textcolor{blue}{\textbf{\therequirement}}%
    \ifx&#1&\else\quad\textbf{#1}\fi  % 可选标题
    \par\nopagebreak
}{%
    %\ifx&#2&\else\label{req:#2}\fi  % 可选标签
    \IfNoValueTF{#2}{}{\label{req:#2}}  % 可选标签
    \par\medskip
}


% counters:
\newcounter{thmcounter}[subsection]
\setcounter{thmcounter}{0}
\newcounter{dfncounter}[subsection]
\setcounter{dfncounter}{0}
\renewcommand{\thethmcounter}{\arabic{section}.\arabic{subsection}.\arabic{thmcounter}}
\renewcommand{\thedfncounter}{\arabic{section}.\arabic{subsection}.\arabic{dfncounter}}

% definition:
\NewDocumentEnvironment{dfn}{s o m s o O{contributor}}{\refstepcounter{dfncounter}\begin{tcolorbox}
    [enhanced, colback=white, boxrule=0pt, frame hidden, borderline west={0.7mm}{0.1mm}{blue}, breakable, after upper={\IfBooleanTF{#4}{\hfill\textcolor{codegray}{\textit{by~#6}}}{}}]
    {\textbf{{\large Definition \thedfncounter :~#3}\IfNoValueTF{#5}{}{\IfBlankTF{#5}{}{(#5)}}}}
    \IfBooleanTF{#1}{}{\newline}
}{\IfNoValueTF{#2}{}{\label{dfn:#2}}\end{tcolorbox}}

% theorem:
\NewDocumentEnvironment{thm}{s o m s o O{contributor}}{\refstepcounter{thmcounter}\begin{tcolorbox} % bool(newline), label, name, bool(contributor displayed), subname, contributor
    [enhanced, boxrule=0pt, frame hidden, borderline west={0.7mm}{0.1mm}{blue}, colback=blue, breakable, after upper={\IfBooleanTF{#4}{\hfill\textcolor{codegray}{\textit{by~#6}}}{}}]
    {\textbf{{\large Theorem \thethmcounter :~#3}\IfNoValueTF{#5}{}{\IfBlankTF{#5}{}{(#5)}}}}
    \IfBooleanTF{#1}{}{\newline}
}{\IfNoValueTF{#2}{}{\label{thm:#2}}\end{tcolorbox}}

\begin{document}

\section{系统需求}
\subsection{安全需求}

\begin{requirement}[用户认证][auth]
系统应支持多因素认证。
\end{requirement}

\begin{requirement}[数据加密][encrypt]
所有传输数据必须加密。
\end{requirement}

\hyperref[req:auth]{认证需求} 和 \hyperref[req:encrypt]{加密需求} 是系统的核心安全需求。

如需求 \ref{req:auth} 和 \ref{req:encrypt} 所示...

\subsection{Test}

\begin{dfn}
    [self-adjoint]
    {自伴随算子}
    [Self-adjoint Operator]
    设 $H$ 是 Hilbert 空间, $T: H \to H$ 是一个有界线性算子. 定义 $T$ 是一个自伴算子当且仅当 \(T=T^*\).
\end{dfn}

\begin{dfn}
    [normal]
    {正规算子}
    [Normal Operator]
    设 $H$ 是 Hilbert 空间, $T: H \to H$ 是一个有界线性算子. 定义 $T$ 是一个正规算子当且仅当 \(TT^*=T^*T\).
\end{dfn}

\begin{dfn}
    [unitary]
    {幺正算子}
    [Unitary Operator]
    设 $H$ 是 Hilbert 空间, $T: H \to H$ 是一个有界线性算子. 定义 $T$ 是一个幺正算子当且仅当 \(T^*T=TT^*=\mathrm{Id}\).
\end{dfn}

\begin{thm}
    [spectral]
    {正规算子的谱定理}
    [Spectral Theorem for Normal Operators]
    设 $H$ 是 Hilbert 空间, $T: H \to H$ 是一个正规算子. 则存在一个唯一的投影测度 $E$ 定义在 $T$ 的谱上, 使得 \(T=\int_{\sigma(T)} \lambda dE(\lambda)\).
\end{thm}

\ref{dfn:self-adjoint} 定义了自伴随算子. \\
\ref{thm:spectral} 给出了正规算子的谱定理. \\
\ref{dfn:normal} 定义了正规算子. \\
\ref{dfn:unitary} 定义了幺正算子. \\

\end{document}